\documentclass[conference,onecolumn]{IEEEtran}
\IEEEoverridecommandlockouts
% Single-column review draft based on the supplied ICAPPS/IEEE template.
% Remove \IEEEoverridecommandlockouts if no funding footnote is added.

\usepackage{cite}
\usepackage{amsmath,amssymb,amsfonts}
\usepackage{algorithmic}
\usepackage{graphicx}
\usepackage{textcomp}
\usepackage{xcolor}
\usepackage{booktabs}
\usepackage{tabularx}
\usepackage{array}
\usepackage{url}

\def\BibTeX{{\rm B\kern-.05em{\sc i\kern-.025em b}\kern-.08em
    T\kern-.1667em\lower.7ex\hbox{E}\kern-.125emX}}
\newcolumntype{Y}{>{\raggedright\arraybackslash}X}

\begin{document}

\title{Has Speaker Diarization Actually Improved? Evaluating Five Years of Model Progress Beyond Standard Adult Speech}

% Replace the placeholder names, affiliations, and email addresses before
% external circulation or submission. Add or remove author blocks as needed.
\author{
\IEEEauthorblockN{First Author}
\IEEEauthorblockA{\textit{Department or School} \\
\textit{Institution Name}\\
City, Country \\
email or ORCID}
\and
\IEEEauthorblockN{Second Author}
\IEEEauthorblockA{\textit{Department or School} \\
\textit{Institution Name}\\
City, Country \\
email or ORCID}
}

\maketitle

\begin{abstract}
Speaker diarization has progressed from embedding--clustering pipelines to
overlap-aware neural systems, end-to-end speaker-activity models, and
generative speech models that jointly emit speaker labels, timestamps, and
words. However, improvements on established adult-speech benchmarks do not
show whether this progress transfers to underrepresented conversational
settings. This paper proposes a common evaluation of four architectural
families on conventional adult meetings, African-accented medical dialogue,
adult--child interaction, and English-primary code-switched conversation. It
asks whether architectural progress produces gains across domains, whether
those gains close transfer gaps, and which diarization failures have or have
not improved. Primary evaluation uses time-based diarization measures;
text generated by unified models is not evaluated as an automatic speech
recognition outcome. The intended contribution is an architecture-aware
robustness audit connecting model evolution to overlap, boundary, short-turn,
speaker-counting, and speaker-confusion failures.
\end{abstract}

\begin{IEEEkeywords}
speaker diarization, African-accented speech, adult--child interaction,
code-switching, speech language models, domain robustness
\end{IEEEkeywords}

\section{Introduction}

Speaker diarization recovers \emph{who spoke when} by assigning anonymous
speaker labels to time intervals. Recent systems range from modular
voice-activity-detection (VAD), embedding, and clustering pipelines to
end-to-end and generative architectures \cite{park2022review}. Yet published
improvements are commonly established on conventional adult-speech benchmarks,
leaving uncertain whether progress transfers to conversational settings that
differ in accent, age composition, domain, overlap, and turn structure.

This study will audit four architectural families under one protocol. The task
is anonymous, within-recording speaker attribution, not identification of a
person across recordings. Automatic speech recognition (ASR) is outside the
primary research questions. When a generative system also returns words, only
its speaker labels and timing are used for primary diarization scoring.

\section{Datasets}

Table~\ref{tab:datasets} gives the final evaluation view. AMI is a conventional
adult reference rather than a matched causal control. AfriSpeech is restricted
to its medical subset. Playlogue is retained because its dyadic structure
closely matches AfriSpeech and most Bangor sessions. Bangor contains
English-primary conversations with natural English--Spanish switching.

\begin{table}[t]
\caption{Final Evaluation Dataset View}
\label{tab:datasets}
\centering
\footnotesize
\setlength{\tabcolsep}{3pt}
\begin{tabularx}{\columnwidth}{@{}p{0.18\columnwidth}p{0.18\columnwidth}p{0.25\columnwidth}Y@{}}
\toprule
\textbf{Dataset} & \textbf{Subset} & \textbf{Speakers} & \textbf{Role and annotation coverage} \\
\midrule
AMI Mix-Headset \cite{carletta2005ami} & 16 sessions; 9.06 h & One 3-speaker;
fifteen 4-speaker & Adult reference; 63.78 min annotated overlap. \\
AfriSpeech-Dialog \cite{sanni2025afri} & Medical only; 17 sessions; 1.77 h &
All 2-speaker & African-accented medical transfer; overlap insufficient. \\
Playlogue v1 \cite{kalanadhabhatta2024playlogue} & Official test; 34 sessions;
8.11 h & All 2-speaker; 2.25 h child speech & Adult--child transfer; released
RTTM has no reference overlap. \\
Bangor Miami \cite{deuchar2014bangor} & 28 English-primary sessions; 14.85 h &
Twenty-six 2-, one 3-, and one 4-speaker & Code-switch transfer; 98.48 min
annotated overlap. \\
\bottomrule
\end{tabularx}
\end{table}

The suite contains 95 sessions and 33.79 audio hours. Across the three target
datasets, 77 of 79 sessions (97.5\%) are dyadic, reducing speaker-count
confounding. Playlogue contains 2.25 hours of child speech (40.71\% of
annotated speech); child and adult errors will be reported separately.

\section{Research Questions and Measures}

Table~\ref{tab:rqs} distinguishes cross-family progress, closure of the
cross-domain performance gap, and the speech conditions in which progress and
transfer are most evident.

\begin{table}[t]
\caption{Research Questions, Purposes, and Measures}
\label{tab:rqs}
\centering
\footnotesize
\setlength{\tabcolsep}{3pt}
\begin{tabularx}{\columnwidth}{@{}p{0.31\columnwidth}p{0.22\columnwidth}Y@{}}
\toprule
\textbf{Research Question} & \textbf{Purpose} & \textbf{Measures} \\
\midrule
\textbf{RQ1.} Have successive architectural shifts in speaker diarization led
to improved accuracy across domains? & Conduct an overall and
within-domain group comparison of the four architectural families. & DER,
JER, missed speech, false alarm, speaker confusion, speaker-count accuracy,
family main effect, planned between-family contrasts, relative error
reduction, and clustered confidence intervals. \\
\textbf{RQ2.} Have these gains closed performance gaps for unconventional and
less represented conversational domains? & Test transfer and
whether the target-domain penalty becomes smaller in later architectural
families. & Target-minus-adult DER/JER gap, change in that gap relative to G1,
family-by-domain interaction, relative gap reduction, rank correlation and
rank reversals, with clustered confidence intervals. \\
\textbf{RQ3.} In which speech conditions have gains and cross-domain transfer
been most effective? & Locate improvements and persistent
failure across overlap, boundaries, short turns, speaker counting, and speaker
confusion. & Condition-stratified DER/JER and confusion, overlap DER, boundary
precision/recall/F1, onset/offset error, speaker-count accuracy, and
family-by-condition interaction. \\
\bottomrule
\end{tabularx}
\end{table}

RQ1 tests the overall pattern of progress through model-family group
comparisons. RQ2 tests whether that progress reduces the gap between the adult
reference benchmark and the target settings. RQ3 identifies where gains occur
and failures remain. Overlap-specific claims are restricted to AMI and Bangor;
Playlogue supports child/adult-conditioned, boundary, short-turn, counting, and
confusion analyses.

\section{Literature-Grounded Architecture Taxonomy}

The comparison strata in Table~\ref{tab:models} are grounded in published
reviews rather than inferred from model release dates. Park \emph{et al.}
organize diarization systems by whether they optimize a diarization objective
and whether components are trained independently or jointly, tracing the move
from traditional modular pipelines through neuralized components and joint
models to fully end-to-end diarization \cite{park2022review}. Serafini
\emph{et al.} independently compare clustering-based, end-to-end neural
diarization (EEND), and speech-separation-guided paradigms, including
overlap-aware clustering variants \cite{serafini2023experimental}. We
operationalize the shared architectural distinctions as G1--G3. G4 is an
explicit post-review extension for recent unified generative systems; it is
not presented as a category proposed by either review. Accordingly, ``G1'' to
``G4'' are study strata describing increasing integration, not universally
accepted or perfectly chronological eras. One model will represent each
family; the entries below remain candidates until the balanced pilot confirms
eligibility. Checkpoint versions and release dates will be fixed and reported.

\begin{table}[t]
\caption{Architectural Families, Shifts, and Representative Models}
\label{tab:models}
\centering
\footnotesize
\setlength{\tabcolsep}{3pt}
\begin{tabularx}{\columnwidth}{@{}p{0.18\columnwidth}p{0.48\columnwidth}Y@{}}
\toprule
\textbf{Family} & \textbf{Architectural Shift} & \textbf{Representative Model} \\
\midrule
\textbf{G1: Embedding--clustering cascade} & Explicit VAD/segmentation,
speaker-embedding extraction, and global clustering; speaker consistency is
determined by embedding similarity. & NeMo Clustering Diarizer: MarbleNet VAD,
TitaNet-Large embeddings, and spectral clustering \cite{koluguri2022titanet}. \\
\textbf{G2: Neuralized overlap-aware modular} & Learned powerset speaker
activity improves local segmentation and overlap handling, while embeddings
and clustering preserve recording-level speaker consistency. & pyannote
Community-1 \cite{plaquet2023powerset,bredin2023pyannote}. \\
\textbf{G3: End-to-end discriminative} & A sequence model directly predicts
frame-level speaker activity without a separate embedding--clustering backend.
& NVIDIA Streaming Sortformer 4-Speaker v2.1 \cite{park2024sortformer}. \\
\textbf{G4: Unified generative} & A long-context model generates timestamps,
anonymous speaker tokens, and words as one structured sequence, shifting
diarization from frame classification to token generation. & Microsoft
VibeVoice-ASR \cite{peng2026vibevoice}. \\
\bottomrule
\end{tabularx}
\end{table}

Model selection will follow a preregistered, eligibility-first rule. A candidate
must have a public and versionable checkpoint, documented architecture,
reproducible inference, time-aligned anonymous speaker output, and compatibility
with the common audio and scoring protocol. Among eligible candidates within a
family, downloads measured on the same platform over a fixed 30-day window will
serve only as an adoption tie-breaker; the snapshot date and counts will be
reported. This prevents repository age, access gating, automatic package
downloads, or compute requirements from becoming the primary scientific
selection criterion.

Closed models can provide external sensitivity comparisons but cannot support
claims about undisclosed internal architectures. Neural-codec dialogue systems
such as Moshi \cite{defossez2024moshi} remain related work: Moshi models
predefined user and system streams rather than inferring anonymous speakers
from mixed multiparty audio. Neural-codec tokenization is an adjacent
speech-language-model development, not by itself a diarization generation.

\section{Experimental Protocol}

All systems will receive the same versioned manifest and standardized 16-kHz
mono recordings. The primary condition will be fully automatic. A secondary
oracle-speaker-count condition will separate counting errors from
speaker-assignment errors. Standalone outputs and parsed generative outputs
will be converted to RTTM. Complete recordings will be used when supported;
any required chunking and cross-window speaker stitching will be fixed after
the balanced pilot and reported as a model constraint. Malformed output,
missing time coverage, and inability to represent simultaneous speakers will
be reported rather than silently repaired.

Primary scoring will use zero collar with overlapped speech included. A
0.25-second-collar sensitivity analysis will support comparison with prior
work. Results will include recording- or participant-clustered bootstrap
confidence intervals and recording-level macro averages. Cross-dataset analyses will adjust for overlap
proportion, turn rate, median turn duration, speech proportion, speaker
imbalance, signal-to-noise ratio, bandwidth, and window duration. Because
accent, age, domain, and recording conditions are not independently
manipulated, conclusions will concern transfer and observed performance gaps,
not isolated causal effects of accent or age. A balanced pilot will validate
the annotations, parsers, scoring protocol, model limits, cost, and data
governance before full evaluation.

\section{Expected Contribution and Novelty}

The contribution is a cross-generational, architecture-aware robustness audit,
not a new diarization architecture. Existing AfriSpeech-Dialog and Playlogue
studies benchmark particular systems but do not establish whether improvements
across model families transfer consistently or mainly raise standard-benchmark
averages. The paper will contribute: (1) one evaluation protocol spanning four
architectural families and four contrasting settings; (2) evidence of whether
generational gains narrow or preserve performance gaps in African-accented
medical, adult--child, and English-primary code-switched speech; (3)
condition-level links between architectural shifts and overlap, boundary,
short-turn, counting, and confusion errors; and (4) a reproducible distinction
among open modular, end-to-end, and generative diarization systems. The central
claim is whether five years of architectural progress have produced reliable
speaker attribution beyond conventional evaluation settings.

\newpage

\section{Dataset Progress and Experiment-to-Paper Roadmap}

All source audio remains unchanged. The reproducible preparation pipeline has
standardized audio to 16-kHz mono, generated or collected RTTM and
Un-partitioned Evaluation Map (UEM) references, and validated paths, durations,
and reference boundaries with zero structural errors. The prepared superset
contains 124 recordings and 38.81 hours; the final medical-only evaluation view
contains 95 recordings and 33.79 hours.

\begin{table}[t]
\caption{Dataset Preparation Status and Remaining Reference Actions}
\label{tab:progress}
\centering
\footnotesize
\setlength{\tabcolsep}{3pt}
\begin{tabularx}{\columnwidth}{@{}p{0.13\columnwidth}p{0.43\columnwidth}Y@{}}
\toprule
\textbf{Dataset} & \textbf{Completed} & \textbf{Remaining action} \\
\midrule
AMI & Official test audio, RTTM, and UEM standardized and validated. & Freeze
checksums and comparison-paper scoring conventions. \\
AfriSpeech & Forty-six usable conversations prepared; final view selects 17
medical sessions. & Export medical-only manifest and manually audit filtered
nonpositive timestamp pairs. \\
Playlogue & Official test audio trimmed, resampled, and paired with RTTM/UEM. &
Document zero overlap in all released RTTMs and validate child/adult scoring. \\
Bangor & Twenty-eight English-primary sessions converted from CHAT to RTTM/UEM;
code-switch metadata retained. & Align 35 untimed trusted tiers or mask their
surrounding intervals; retain all sessions. \\
\bottomrule
\end{tabularx}
\end{table}

\subsection{Required Experiments}

\begin{enumerate}
\item \textbf{Lock datasets and references:} resolve Bangor timing, export the
medical-only AfriSpeech view, add checksums and version identifiers, and freeze
the metric-eligibility matrix.
\item \textbf{Finalize model selection:} verify one candidate per family;
freeze checkpoint/API versions, licenses, context limits, compute needs,
inference settings, and output parsers.
\item \textbf{Balanced pilot:} sample every dataset and speaker-count regime to
test output parsing, overlap support, speaker-count behavior, determinism,
runtime, memory, cost, and the full-recording versus common-window policy.
\item \textbf{Primary inference:} run every model on every recording with
automatic speaker count while preserving raw output, parsed RTTM, logs,
runtime, peak memory, failures, and configuration hashes.
\item \textbf{Controlled inference:} repeat with oracle speaker count where
supported; run the preregistered chunk/stitch condition only if model context
limits require it.
\item \textbf{Overall scoring:} compute DER, JER, missed speech, false alarm,
speaker confusion, speaker-count error, malformed-output rate, and failure rate
under zero and 0.25-second collars.
\item \textbf{Condition scoring:} analyze short turns, boundaries, rapid
speaker changes, imbalance, and per-speaker confusion; overlap on AMI/Bangor;
child/adult errors on Playlogue; and switched/non-switched turns on Bangor.
\item \textbf{Statistical analysis:} use recording-level macro averages and
clustered bootstrap intervals; test family, dataset, family-by-dataset, and
family-by-condition effects; report corrected planned contrasts, relative
error and gap reduction, rank correlation, and rank reversals.
\item \textbf{Robustness and error audit:} repeat both collar settings; test
with and without incidental or highly imbalanced speakers; manually inspect a
stratified failure sample; distinguish parser from acoustic-model failures.
\end{enumerate}

\subsection{Final-Paper Analyses and Completion Gates}

The four-page paper requires: Table I for RQs and measures; Table II for the
architecture taxonomy and selected checkpoints; a compact dataset table; a
family-by-dataset DER/JER results table with transfer gaps and confidence
intervals; one architectural-progress figure; and one compact condition-gain
heatmap with ineligible cells explicitly marked. The reproducibility package
will include locked manifests, lawful preparation instructions, RTTM/UEM,
parsers, scoring configuration, model versions, environment lockfile, seeds,
and raw-to-result provenance.

Full paper writing begins after four gates: (G1) references and metric
eligibility are frozen; (G2) all four parsers pass the pilot; (G3) the complete
model--recording matrix or explicit failure log exists; and (G4) RQ1--RQ3
tables, confidence intervals, and sensitivity analyses are generated. Writing
will proceed in the order Methods, Results, Discussion/limitations, Related
Work/taxonomy justification, then Abstract and Introduction after the central
empirical claim is known.

\begin{thebibliography}{00}

\bibitem{park2022review}
T. J. Park, N. Kanda, D. Dimitriadis, K. J. Han, S. Watanabe, and S. Narayanan,
``A review of speaker diarization: Recent advances with deep learning,''
\emph{Computer Speech \& Language}, vol. 72, Art. no. 101317, 2022,
doi: 10.1016/j.csl.2021.101317.

\bibitem{serafini2023experimental}
L. Serafini, S. Cornell, G. Morrone, E. Zovato, A. Brutti, and S. Squartini,
``An experimental review of speaker diarization methods with application to
two-speaker conversational telephone speech recordings,'' \emph{Computer
Speech \& Language}, vol. 82, Art. no. 101534, 2023,
doi: 10.1016/j.csl.2023.101534.

\bibitem{carletta2005ami}
J. Carletta \emph{et al.}, ``The AMI meeting corpus: A pre-announcement,'' in
\emph{Proc. MLMI}, 2005, pp. 28--39.

\bibitem{sanni2025afri}
M. Sanni \emph{et al.}, ``Afrispeech-Dialog: A benchmark dataset for
spontaneous English conversations in healthcare and beyond,'' in \emph{Proc.
NAACL-HLT}, 2025, pp. 8399--8417.

\bibitem{kalanadhabhatta2024playlogue}
M. Kalanadhabhatta, M. M. Rastikerdar, T. Rahman, A. S. Grabell, and D. Ganesan,
``Playlogue: Dataset and benchmarks for analyzing adult-child conversations
during play,'' \emph{Proc. ACM Interact. Mob. Wearable Ubiquitous Technol.},
vol. 8, no. 4, 2024, doi: 10.1145/3699775.

\bibitem{deuchar2014bangor}
M. Deuchar, P. Davies, J. R. Herring, M. C. Parafita Couto, and D. Carter,
``Building bilingual corpora,'' in \emph{Advances in the Study of
Bilingualism}, E. M. Thomas and I. Mennen, Eds. Bristol, U.K.: Multilingual
Matters, 2014, pp. 93--111.

\bibitem{koluguri2022titanet}
N. R. Koluguri, T. Park, and B. Ginsburg, ``TitaNet: Neural model for speaker
representation with 1D depth-wise separable convolutions and global context,''
in \emph{Proc. IEEE ICASSP}, 2022, pp. 8102--8106.

\bibitem{bredin2023pyannote}
H. Bredin, ``pyannote.audio 2.1 speaker diarization pipeline: Principle,
benchmark, and recipe,'' in \emph{Proc. Interspeech}, 2023.

\bibitem{plaquet2023powerset}
A. Plaquet and H. Bredin, ``Powerset multi-class cross entropy loss for neural
speaker diarization,'' in \emph{Proc. Interspeech}, 2023, pp. 3222--3226.

\bibitem{park2024sortformer}
T. Park \emph{et al.}, ``Sortformer: Seamless integration of speaker
diarization and ASR by bridging timestamps and tokens,'' arXiv:2409.06656,
2024.

\bibitem{peng2026vibevoice}
Z. Peng \emph{et al.}, ``VIBEVOICE-ASR technical report,'' arXiv:2601.18184,
2026.

\bibitem{defossez2024moshi}
A. D\'efossez, L. Mazar\'e, M. Orsini, A. Royer, P. P\'erez, H. J\'egou,
E. Grave, and N. Zeghidour, ``Moshi: A speech-text foundation model for
real-time dialogue,'' arXiv:2410.00037, 2024.

\end{thebibliography}

\end{document}
